
\section{Problem Formulation}

\subsection{Contexts, policies and the two-level decision}

Let $\mathcal{S}$ be a set of operating contexts and $\mathcal{A}$ a finite set of
retained routing policies. A context $s\in\mathcal{S}$ contains only quantities known
\emph{before} the policy is activated; here it is the triple (scheduled primary demand,
planned restriction configuration, signal-control plan), and no realised travel time,
queue, flow or emission is part of it. A policy $a\in\mathcal{A}$ maps a route request to
one path from a finite admissible catalogue according to its own objective.

The decision has two levels: \emph{Level 1}, policy selection, maps the context $s$ to one
policy $a\in\mathcal{A}$; \emph{Level 2}, route assignment, is performed by that policy
according to its own objective. Only Level 1 is learned. The unit of analysis is the context, not the vehicle and not the
replication: a context is one instance in the algorithm-selection sense, all of its
replications move together through every evaluation protocol, and the selector never
observes the realised outcome of any candidate policy in the context it decides about.

\subsection{Decision cost, regret and headroom}

Let $y_{sak}$ be the mean outcome of criterion $k$ in context $s$ under policy $a$,
averaged over replications after run-level aggregation. With positive fixed references
$r_k$ and weights $w_k\ge0$, $\sum_k w_k=1$, the scalar decision cost is
\begin{equation}
C_{sa}=\sum_k w_k\,y_{sak}/r_k ,
\label{eq:cost}
\end{equation}
lower being preferred. The references are fixed for the whole benchmark rather than
recomputed per context, so costs are comparable across contexts and a context in which
every policy is slow is not rescaled into apparent equality; one cost unit is one
benchmark-average criterion value.

The best policy actually observed in a context is
$a_s^{*}\in\arg\min_{a\in\mathcal{A}} C_{sa}$, with cost $C_{s}^{\mathrm{VB}}$. Selecting
$a_s^{*}$ everywhere defines the \emph{hindsight} (virtual best) policy: a bound, not a
method, since it uses the realised performance of every candidate. For a selection rule
$\pi:\mathcal{S}\rightarrow\mathcal{A}$ the \emph{decision regret}
$R_s(\pi)=C_{s,\pi(s)}-C^{\mathrm{VB}}_{s}\ge0$ measures the cost of the decision, not the
error of a prediction. The strongest deployable non-adaptive rule is the \emph{single best
fixed policy} $a^{\mathrm{SB}}=\arg\min_{a}|\mathcal{S}|^{-1}\sum_{s}R_s(a)$, and the
quantity adaptation can win is the \emph{decision headroom} $\bar H$, of which a rule
recovers the fraction $G(\pi)$:
\begin{equation}
\bar H=\frac{1}{|\mathcal{S}|}\sum_{s\in\mathcal{S}}
\bigl(C_{s,a^{\mathrm{SB}}}-C^{\mathrm{VB}}_{s}\bigr),
\qquad
G(\pi)=\frac{\bar H-\overline{R(\pi)}}{\bar H} .
\label{eq:headroom}
\end{equation}
Thus $G=1$ for the hindsight bound, $G=0$ for the single best fixed policy, and $G<0$ for
a rule worse than doing nothing adaptive. Because $\bar H$ is defined on a finite context
set, $G$ is a benchmark-specific decision statistic. It is \emph{not} a percentage
reduction in travel time, emissions or any other physical quantity; the two are reported
separately throughout.

\subsection{Three properties that must not be confused}

\emph{Complementarity} asks whether different policies win in different contexts. It is a
property of the portfolio and environment alone, established by observing that $a^{*}_{s}$
is not constant over $\mathcal{S}$, and it says nothing about how much the variation is
worth. \emph{Decision value} asks how much cost the variation carries, and is $\bar H$ in
\eqref{eq:headroom}; a portfolio can be strongly complementary while $\bar H$ is
negligible, if the single best fixed policy is nearly optimal wherever it is not optimal,
so complementarity is necessary for decision value but not sufficient.
\emph{Deployability} asks whether the variation can be exploited from pre-decision
information on contexts not used for fitting, and is $G(\pi)$ under a held-out protocol
together with evidence about the range of contexts over which the rule remains valid;
positive decision value does not imply it, and a rule deployable on one region of the
context space need not be deployable outside it. All three are measured separately in
Section~\ref{sec:results}, and this study reports a different answer for each.
